% block: shifted-systems-heading pages:175,176 type:section
\subsection{{Shifted Linear Systems}}\label{sec:shifted-linear-systems}
{The Krylov subspace methods in Chap.\ \ref{chap:krylov-classification} are attractive for solving the following linear systems:}
\begin{equation}\label{eq:shifted-linear-systems}A^{(k)}\bm x^{(k)}=\bm b,\quad k=1,2,\ldots,m,\end{equation}
{where $A^{(k)}$ is a nonsingular matrix of the form}
\begin{equation}\label{eq:shifted-coefficient-matrix}A^{(k)}:=A+\sigma_k I,\end{equation}
{with $I$ being the identity matrix and $\sigma_k\in\mathbb C$. The linear systems \eqref{eq:shifted-linear-systems} are referred to as shifted linear systems. For simplicity, we also use the following notation:}
\[(A+\sigma I)\bm x^{\sigma}=\bm b.\]
{Here $\alpha^{\sigma},\beta^{\sigma},\ldots$ and $\bm a^{\sigma},\bm b^{\sigma},\ldots$ denote scalars and vectors related to the shifted linear system $(A+\sigma I)\bm x^{\sigma}=\bm b$. Note that $\alpha^{\sigma}$ is not the $\alpha$ to the power of $\sigma$.}
% block: shifted-krylov-invariance page:176 type:ordinary-paragraph
{If a Krylov subspace method with the initial guess $\bm x_0=\bm0$ is applied to \eqref{eq:shifted-linear-systems}, then from \eqref{eq:krylov-iterate-condition} the Krylov subspace method finds the approximate solutions of \eqref{eq:shifted-linear-systems} over the following subspaces:}
\begin{equation}\label{eq:shifted-krylov-membership}\bm x_n^{(k)}\in\mathcal K_n(A^{(k)},\bm b),\quad k=1,2,\ldots,m.\end{equation}
{From the definition of Krylov subspace \eqref{eq:krylov-subspace-span}, the following shift invariance property holds true:}
\begin{equation}\label{eq:krylov-shift-invariance}\mathcal K_n(A,\bm b)=\mathcal K_n(A^{(k)},\bm b).\end{equation}
% block: shift-invariance-example page:176 type:ordinary-paragraph
{This can be easily shown by induction and the definition of a Krylov subspace. As an example, for the case $n=3$ it follows that}
\[\begin{aligned}\mathcal K_3(A+\sigma I,\bm b)&=\{c_1\bm b+c_2(A+\sigma I)\bm b+c_3(A+\sigma I)^2\bm b:c_1,c_2,c_3\in\mathbb C\}\\
&=\{(c_1+c_2\sigma+c_3\sigma^2)\bm b+(c_2+2c_3\sigma)A\bm b+c_3A^2\bm b:c_1,c_2,c_3\in\mathbb C\}\\
&=\{d_1\bm b+d_2A\bm b+d_3A^2\bm b:d_1,d_2,d_3\in\mathbb C\}\\
&=\mathcal K_3(A,\bm b).\end{aligned}\]
{In the third equation, the relation between $c_i$ and $d_i$ is given by}
\[\begin{bmatrix}d_1\\d_2\\d_3\end{bmatrix}=\begin{bmatrix}1&\sigma&\sigma^2\\0&1&2\sigma\\0&0&1\end{bmatrix}\begin{bmatrix}c_1\\c_2\\c_3\end{bmatrix}.\]
% block: collinear-shift-invariance page:176 type:ordinary-paragraph
{Since the $3\times3$ matrix is nonsingular, for any $d_1,d_2,d_3$ there exist $c_1,c_2,c_3$. Thus the third equation holds true. From this example, we see that the shift invariance property can be slightly extended as follows:}
\[\mathcal K_n(A,\bm v)=\mathcal K_n(A^{(k)},\bm w),\quad\bm v=c\bm w.\]
{If $\bm v=c\bm w$ holds for some $c$, the two vectors are referred to as collinear.}
% block: shared-shifted-subspace pages:176,177 type:ordinary-paragraph
{Now from the shift invariance property \eqref{eq:krylov-shift-invariance}, it follows that \eqref{eq:shifted-krylov-membership} can be rewritten as}
\[\bm x_n^{(k)}\in\mathcal K_n(A,\bm b),\quad k=1,2,\ldots,m.\]
{This means that only one Krylov subspace, $\mathcal K_n(A,\bm b)$, is required for obtaining the approximate solutions $\bm x_n^{(1)},\bm x_n^{(2)},\ldots,\bm x_n^{(m)}$. Note that if the shift-invariance property does not hold, $m$ different Krylov subspaces are required. In the following sections, we will see how to utilize the shift-invariance property and how this property is incorporated into Krylov subspace methods.}
